\documentclass[11pt,a4paper]{article}
\usepackage[T1]{fontenc}
\usepackage[utf8]{inputenc}
\usepackage[ngerman]{babel}
\usepackage[margin=2.5cm]{geometry}
\usepackage[hidelinks]{hyperref}
\usepackage{array}
\setlength{\parindent}{0pt}
\setlength{\emergencystretch}{2em}
\renewcommand{\arraystretch}{1.25}
\begin{document}
\shorthandoff{"}
\begin{center}
  {\Large\textbf{Webtechnologien -- Übungsblatt 2}}
\end{center}
\section*{Aufgabe 5: Event Listeners and Target}
\textbf{Aufgabe:} Die Konsolenausgabe beim Klick auf den Button vorhersagen
und prüfen, \texttt{target}, \texttt{currentTarget} und \texttt{this}
erklären. Danach \texttt{ex5.js} so ergänzen, dass nur die Informationen
des Buttons ausgegeben werden; \texttt{ex5.html} bleibt unverändert.

\medskip
\textbf{Ausgangscode:} \texttt{ex5.js} hängt je einen \texttt{click}-Listener
an den Button \texttt{\#button} und an das umgebende \texttt{div}
\texttt{\#outerArea}. Beide geben \glqq inner:\grqq{} bzw.
\glqq outer:\grqq{} und danach \texttt{target}, \texttt{currentTarget} und
\texttt{this} aus. Durch die String-Verkettung erscheint jeweils nur der
Objekttyp, z.\,B. \texttt{[object HTMLButtonElement]}.

\medskip
\textbf{5.1 und 5.2 Vorhersage und Ausgabe:} Zuerst erscheinen
\glqq inner:\grqq{} und drei Zeilen, danach \glqq outer:\grqq{} und drei Zeilen:
\begin{center}
\begin{tabular}{|l|l|l|}
\hline
\textbf{Zeile} & \textbf{inner (Listener am Button)} & \textbf{outer (Listener am div)} \\ \hline
\texttt{target} & \texttt{[object HTMLButtonElement]} & \texttt{[object HTMLButtonElement]} \\ \hline
\texttt{currentTarget} & \texttt{[object HTMLButtonElement]} & \texttt{[object HTMLDivElement]} \\ \hline
\texttt{this} & \texttt{[object HTMLButtonElement]} & \texttt{[object HTMLDivElement]} \\ \hline
\end{tabular}
\end{center}
Grund ist Event-Bubbling: Das \texttt{click}-Event läuft vom Button nach oben
durch \texttt{div}, \texttt{body}, \texttt{html} und \texttt{document}.
Jeder Listener auf diesem Weg wird ausgeführt. Die tatsächliche Ausgabe in
Chrome stimmt mit der Vorhersage überein. Ein Klick auf die blaue Fläche
neben dem Button gibt nur \glqq outer:\grqq{} aus, dort sind alle drei
Werte \texttt{[object HTMLDivElement]}.

\medskip
\textbf{5.3 target, currentTarget und this:}
\begin{itemize}
  \item \texttt{target}: das Element, auf dem das Event ausgelöst wurde,
    also das tatsächlich angeklickte, innerste Element. Bleibt beim
    Bubbling gleich.
  \item \texttt{currentTarget}: das Element, an dem der gerade laufende
    Listener hängt. Ändert sich beim Hochwandern des Events.
  \item \texttt{this}: In einer normalen \texttt{function} als Listener
    dasselbe wie \texttt{currentTarget}. Bei einer Pfeilfunktion wäre
    \texttt{this} nicht das Element, sondern kommt aus dem umgebenden
    Kontext.
\end{itemize}

\textbf{5.4 Lösung:} Am Ende von \texttt{ex5.js} (ab Zeile 15) ist ein
zusätzlicher \texttt{click}-Listener auf \texttt{\#button} hinzugefügt,
der \texttt{event.stopPropagation()} aufruft. Das Event wandert dadurch
nicht mehr zum \texttt{div}. Beide Listener am Button laufen weiterhin,
weil \texttt{stopPropagation} nur den Weg zu anderen Elementen stoppt.
Der bestehende Code ist unverändert.

\medskip
\textbf{Hinweis:} Alternativ könnte der äußere Listener mit
\begin{center}
\texttt{if (event.target !== event.currentTarget) return;}
\end{center}
beginnen. Dafür müsste aber der vorhandene Code geändert werden.

\medskip
\textbf{Test:} In Chrome, \texttt{ex5.html} direkt geöffnet.
Klick auf den Button: nur die vier Zeilen von \glqq inner:\grqq.
Klick auf die blaue Fläche: weiterhin \glqq outer:\grqq.

\medskip
{\small \textbf{Hilfsmittel:} Claude Code (Claude Opus 5.5) für Erstellung, Prüfung und Dokumentation.}
\end{document}
